
\section{\BookRecordIndex{T0266}平均值不等式}\label{sec:4-3}
\zhindex{pingjunzhibudengshi}{平均值不等式}

本节将为小能量 \BookRecordIndex{T0003}$J$-全纯曲线的一阶导数建立如下\BookRecordIndex{T0266}平均值不等式。记
\[
  B_r(z):=\{\zeta\in\mathbb C\mid |\zeta-z|\le r\}
\]
为以 $z$ 为中心、半径为 $r$ 的闭圆盘，并简记 $B_r:=B_r(0)$。

\begin{lemma}\label{lem:4-3-1}
设 $(M,J)$ 是紧殆复流形，$L\subset M$ 是紧\BookRecordIndex{T0290}全实子流形。在 $M$ 上选取一个黎曼度量。
那么存在常数 $c,\delta>0$，使得下列结论成立。
\begin{enumerate}[label=\textup{(\roman*)},leftmargin=3em]
\item 若 $r>0$，且 $u:B_r\to M$ 是 \BookRecordIndex{T0003}$J$-全纯曲线，则
\[
  \int_{B_r}|du|^2<\delta
  \qquad\Longrightarrow\qquad
  |du(0)|^2
  \le\frac{c}{r^2}\int_{B_r}|du|^2.
\]
\item 若 $r>0$，且
\[
  u:(B_{2r}\cap\mathbb H,B_{2r}\cap\mathbb R)
  \longrightarrow(M,L)
\]
是 \BookRecordIndex{T0003}$J$-全纯曲线，则
\[
  \int_{B_{2r}\cap\mathbb H}|du|^2<\delta
  \qquad\Longrightarrow\qquad
  \sup_{B_r\cap\mathbb H}|du|^2
  \le\frac{2c}{r^2}
  \int_{B_{2r}\cap\mathbb H}|du|^2.
\]
\end{enumerate}
\end{lemma}

\begin{remark}\label{rem:4-3-2}
把引理~\ref{lem:4-3-1} 中的黎曼度量记为 $g$。该引理的证明表明，常数 $\delta$
依赖于 $M$ 上一个辅助黎曼度量的选取；关于这个辅助度量，$J$ 是\BookRecordIndex{T0350}斜伴随的，
$JTL$ 与 $TL$ 正交，而且 $L$ 是\BookRecordIndex{T0286}全测地的\sourcefootnote{原书此处写作
“and $L$ is totally real”。但 $L$ 为\BookRecordIndex{T0290}全实子流形已是引理~\ref{lem:4-3-1} 的假设，
与度量选取无关；紧接着所引的引理~\ref{lem:4-3-4} 构造的正是使 $L$ \BookRecordIndex{T0286}全测地的度量，
故据此改正。}。称之为 $J$ 与 $L$ 的 \BookRecordIndex{T0018}Frauenfelder 度量，记作
$g_{J,L}$（见下面的引理~\ref{lem:4-3-4}）。把引理~\ref{lem:4-3-1} 在
$g=g_{J,L}$ 时的常数记为 $\delta_{J,L}$ 和 $c_{J,L}$。证明表明
$c_{J,L}=8/\pi$，而 $\delta_{J,L}$ 关于 $J$ 和 $g_{J,L}$ 在 $C^2$ 拓扑下连续。
估计中涉及 $g_{J,L}$ 的 \BookRecordIndex{T0034}Levi--Civita 联络和黎曼曲率张量，以及 $J$ 关于
$g_{J,L}$ 的一阶和二阶协变导数。引理~\ref{lem:4-3-1} 中的常数 $c$ 和 $\delta$
只依赖于 $c_{J,L}$、$\delta_{J,L}$ 以及度量 $g_{J,L}$ 与 $g$ 的比较常数。总之，
当 $L=\varnothing$ 时，引理~\ref{lem:4-3-1} 中的 $\delta$ 关于 $J$ 在 $C^2$
拓扑下连续，而 $c$ 关于 $J$ 在 $C^0$ 拓扑下连续。
\end{remark}

引理~\ref{lem:4-3-1} 的证明依赖\BookRecordIndex{T0251}能量密度
\[
  w:=\frac12|du|^2
\]
所满足的偏微分不等式
\[
  \Delta w\ge-aw^2,
\]
其中
\[
  \Delta:=\partial_s^2+\partial_t^2
\]
是 $\mathbb R^2$ 上的标准\BookRecordIndex{T0220}拉普拉斯算子。每个维数中都有类似结果（例如见
\srcite{336}），但这里我们只考察二元函数。

\begin{lemma}\label{lem:4-3-3}
设 $r>0$ 且 $a\ge0$。若 $C^2$ 函数 $w:B_r\to\mathbb R$ 满足
\begin{equation}\label{eq:4-3-1}
  \Delta w\ge-aw^2,
  \qquad
  w\ge0,
  \qquad
  \int_{B_r}w<\frac{\pi}{8a},
\end{equation}
其中当 $a=0$ 时略去最后一个积分条件
\sourcefootnote{原书在 $a\ge0$ 的同时直接写出 $\int_{B_r}w<\pi/(8a)$，但该商式在
$a=0$ 时无定义；证明的步骤 3 明确把 $a=0$ 作为无须积分小量条件的次调和情形处理，
故按书内论证补明这一约定。}，
则
\begin{equation}\label{eq:4-3-2}
  w(0)\le\frac{8}{\pi r^2}\int_{B_r}w.
\end{equation}
\end{lemma}

\begin{proof}
证明分为五步。

\smallskip
\noindent\textit{步骤 1。}
若 $C^2$ 函数 $w:B_r\to\mathbb R$ 满足
\[
  \Delta w\ge-b,\qquad w\ge0
\]
其中 $b\ge0$ 是常数，则
\[
  w(0)\le\frac{br^2}{8}+\frac{1}{\pi r^2}\int_{B_r}w.
\]
这是\BookRecordIndex{T0079}次调和函数
\[
  v(s,t):=w(s,t)+\frac b4(s^2+t^2)
\]
的\BookRecordIndex{T0266}平均值不等式。为求完整，我们给出证明。由散度定理，
\[
\begin{aligned}
  0
  &\le\frac1\rho\int_{B_\rho}\Delta v
   =\frac1\rho\int_{\partial B_\rho}\frac{\partial v}{\partial\nu}
   =\int_0^{2\pi}\frac d{d\rho}v(\rho e^{i\theta})\,d\theta \\
  &=\frac d{d\rho}
  \left(\frac1\rho\int_{\partial B_\rho}v\right).
\end{aligned}
\]
因此，当 $0<\sigma\le\rho$ 时，
\[
  \frac{1}{2\pi\sigma}\int_{\partial B_\sigma}v
  \le
  \frac{1}{2\pi\rho}\int_{\partial B_\rho}v.
\]
令 $\sigma$ 趋于零，左边趋于 $v(0)$，故
\[
  2\pi\rho\,v(0)\le\int_{\partial B_\rho}v.
\]
从 $0$ 到 $r$ 积分，得到 $v$ 的\BookRecordIndex{T0266}平均值不等式。于是
\[
  w(0)=v(0)
  \le\frac{1}{\pi r^2}\int_{B_r}v
  =\frac{br^2}{8}+\frac{1}{\pi r^2}\int_{B_r}w.
\]
步骤 1 得证。

\smallskip
\noindent\textit{步骤 2。只须对 $r=1$ 证明引理。}
设 $w:B_r\to\mathbb R$ 满足 \eqref{eq:4-3-1}，并定义
$\widetilde w:B_1\to\mathbb R$ 和 $\widetilde a\in\mathbb R$ 为
\[
  \widetilde w(s,t):=w(rs,rt),
  \qquad
  \widetilde a:=ar^2.
\]
那么
\[
  \Delta\widetilde w\ge-\widetilde a\,\widetilde w^2,
  \qquad
  \int_{B_1}\widetilde w
  =\frac1{r^2}\int_{B_r}w
  <\frac{\pi}{8ar^2}
  =\frac{\pi}{8\widetilde a}.
\]
因此，假设引理对 $r=1$ 成立，就有
\[
  w(0)=\widetilde w(0)
  \le\frac8\pi\int_{B_1}\widetilde w
  =\frac{8}{\pi r^2}\int_{B_r}w.
\]
步骤 2 得证。

\smallskip
\noindent\textit{步骤 3。只须对 $a=1$ 证明引理。}
当 $a=0$ 时，结论由步骤 1 得到。因此设 $a>0$，并设
$w:B_r\to\mathbb R$ 满足 \eqref{eq:4-3-1}。那么
$\widetilde w:=aw$ 满足把 \eqref{eq:4-3-1} 中 $a$ 换成 $1$ 后的条件。若结论对
$a=1$ 成立，则 $\widetilde w$ 满足 \eqref{eq:4-3-2}，从而 $w$ 也满足。

\smallskip
\noindent\textit{步骤 4（\BookRecordIndex{T0025}Heinz 技巧）。}
设 $a=r=1$，并定义 $f:[0,1]\to\mathbb R$ 为
\[
  f(\rho):=(1-\rho)^2\sup_{B_\rho}w,
  \qquad
  0\le\rho\le1.
\]
由于 $f(1)=0$ 且 $f$ 非负，存在 $\rho^*\in[0,1)$ 和 $z^*\in B_{\rho^*}$，使得
\[
  f(\rho^*)=\max_{0\le\rho\le1}f(\rho),
  \qquad
  c:=w(z^*)=\sup_{B_{\rho^*}}w.
\]
记
\[
  \varepsilon:=\frac{1-\rho^*}{2}.
\]
那么，在范围
$0<\rho\le\varepsilon$ 内\sourcefootnote{原书此处写作
“$0\le\rho\le\varepsilon$”；但紧随其后的不等式含有 $1/\rho^2$，在 $\rho=0$
处没有定义，故将左端点改为严格不等号。}，
\begin{equation}\label{eq:4-3-3}
  c\le2c^2\rho^2+\frac{1}{\pi\rho^2}\int_{B_1}w.
\end{equation}
为证明这一点，先注意
\[
\begin{aligned}
  \sup_{B_\varepsilon(z^*)}w
  &\le\sup_{B_{\rho^*+\varepsilon}}w
   =\frac{f(\rho^*+\varepsilon)}
          {(1-\rho^*-\varepsilon)^2} \\
  &=\frac{4f(\rho^*+\varepsilon)}{(1-\rho^*)^2}
   \le\frac{4f(\rho^*)}{(1-\rho^*)^2}
   =4c.
\end{aligned}
\]
所以在 $B_\varepsilon(z^*)$ 中，
$\Delta w\ge-w^2\ge-16c^2$。在步骤 1 中取 $b:=16c^2$，并取
$r:=\rho\le\varepsilon$，便得到 \eqref{eq:4-3-3}。步骤 4 得证。

\smallskip
\noindent\textit{步骤 5。引理对 $r=1$ 和 $a=1$ 成立。}
设 $a=r=1$，并沿用步骤 4 的 $\rho^*,z^*,c,\varepsilon$。若
$4c\varepsilon^2\ge1$，则在 \eqref{eq:4-3-3} 中取
\[
  \rho:=\sqrt{\frac1{4c}}\le\varepsilon,
\]
得到
\[
  c=w(z^*)
  \le\frac c2+\frac{4c}{\pi}\int_{B_1}w,
\]
从而 $\int_{B_1}w\ge\pi/8$，这与假设矛盾。因此
$4c\varepsilon^2<1$。再在 \eqref{eq:4-3-3} 中取 $\rho=\varepsilon$，得到
\[
  c
  \le2c^2\varepsilon^2
    +\frac{1}{\pi\varepsilon^2}\int_{B_1}w
  \le\frac c2
    +\frac{1}{\pi\varepsilon^2}\int_{B_1}w.
\]
所以
\[
  w(0)=f(0)
  \le f(\rho^*)
  =(1-\rho^*)^2c
  =4c\varepsilon^2
  \le\frac8\pi\int_{B_1}w.
\]
这就证明了步骤 5 和引理~\ref{lem:4-3-3}。
\end{proof}

下一个引理由 Urs Frauenfelder 在其 Diplom 学位论文 \srcite{120} 中证明。它断言，
存在一个度量，使 $(M,J)$ 的\BookRecordIndex{T0290}全实子流形 $L$ 成为\BookRecordIndex{T0286}全测地子流形。

\begin{lemma}[Frauenfelder]\label{lem:4-3-4}
设 $(M,J)$ 是殆复流形，$L\subset M$ 是满足
$2\dim L=\dim M$ 的\BookRecordIndex{T0290}全实子流形。那么 $M$ 上存在黎曼度量
$g=\langle\,\cdot\,,\,\cdot\,\rangle$，使得：
\begin{enumerate}[label=\textup{(\roman*)},leftmargin=3em]
\item 对 $p\in M$ 以及 $v,w\in T_pM$，
\[
  \langle J(p)v,J(p)w\rangle=\langle v,w\rangle;
\]
\item 对每个 $p\in L$，$J(p)T_pL$ 是 $T_pL$ 的正交补；
\item 关于 $g$，$L$ 是\BookRecordIndex{T0286}全测地子流形。
\end{enumerate}
\end{lemma}

\begin{proof}
在 $L$ 上选取坐标 $x_1,\ldots,x_n$，并把它们延拓为 $M$ 上的坐标
\[
  x_1,\ldots,x_n,y_1,\ldots,y_n,
\]
使得在 $L$ 上
\[
  J\frac{\partial}{\partial x_i}
  =\frac{\partial}{\partial y_i},
  \qquad
  i=1,\ldots,n.
\]
在这些坐标中，把一个度量和\BookRecordIndex{T0088}殆复结构写成
\[
  g(x,y)
  =
  \begin{pmatrix}
    a(x,y) & b(x,y)^T\\
    b(x,y) & c(x,y)
  \end{pmatrix},
  \qquad
  J(x,y)
  =
  \begin{pmatrix}
    A(x,y) & B(x,y)\\
    C(x,y) & D(x,y)
  \end{pmatrix},
\]
其中
\[
  a(x,y)=a(x,y)^T,
  \qquad
  c(x,y)=c(x,y)^T,
\]
而 $b(x,y),A(x,y),B(x,y),C(x,y),D(x,y)$ 都是实 $n\times n$ 矩阵。于是
\[
  B(x,0)=-\boldsymbol 1,
  \qquad
  C(x,0)=\boldsymbol 1,
  \qquad
  A(x,0)=D(x,0)=0.
\]
度量 $g$ 满足 \textup{(i)}、\textup{(ii)}、\textup{(iii)}，当且仅当
\begin{equation}\label{eq:4-3-4}
  a(x,0)=c(x,0),
  \qquad
  b(x,0)=0,
  \qquad
  \partial_{n+i}a(x,0)=0,
  \qquad
  J^TgJ=g.
\end{equation}
为证明这种度量的局部存在性，先选取一个满足 \eqref{eq:4-3-4} 中前三个恒等式以及
\[
  \partial_{n+i}c(x,0)
  +\partial_{n+i}C(x,0)^Tc(x,0)
  +c(x,0)\partial_{n+i}C(x,0)
  =0
\]
的度量 $g$，然后以 $J^TgJ+g$ 代替 $g$。满足 \eqref{eq:4-3-4} 的度量集合在凸组合
下不变；它在乘以满足
\[
  \partial_{n+i}\beta(x,0)=0
\]
的\BookRecordIndex{T0188}截断函数 $\beta=\beta(x,y)$ 时也保持这些恒等式。对\BookRecordIndex{T0188}截断函数的这一条件是内蕴的，
它断言
\begin{equation}\label{eq:4-3-5}
  q\in L,\quad v\in T_qL
  \qquad\Longrightarrow\qquad
  d\beta(q)J(q)v=0.
\end{equation}
因此，选取满足 \eqref{eq:4-3-4} 的局部度量，再用满足 \eqref{eq:4-3-5} 的局部有限
\BookRecordIndex{T0188}截断函数族所成的单位分解进行拼接，即得结论\sourcefootnote{原书此处写作“由有限多个
\BookRecordIndex{T0188}截断函数组成的单位分解”。引理并未假设 $M$ 或 $L$ 紧；一般情形所需且足够的是局部有限
单位分解，故据此改正。}。
\end{proof}

\begin{proof}[引理~\ref{lem:4-3-1} 的证明]
假设 $M$ 上的度量如引理~\ref{lem:4-3-4}。在这一假设下，我们以常数 $c=8/\pi$
证明该引理。对一般度量，由于 $M$ 紧，从而任意两个度量可比较，结论便对某个适当常数
$c$ 成立。

用 $\nabla$ 表示引理~\ref{lem:4-3-4} 中度量的 \BookRecordIndex{T0034}Levi--Civita 联络，并用
\[
  R\in\Omega^2(M,\End(TM))
\]
表示 $\nabla$ 的曲率张量。简记
\[
  \xi:=\partial_su,
  \qquad
  \eta:=\partial_tu,
  \qquad
  w:=\frac12|du|^2:B_r\longrightarrow\mathbb R.
\]
于是
\[
  \xi+J\eta=0,
  \qquad
  w=|\xi|^2=|\eta|^2,
\]
并且
\begin{equation}\label{eq:4-3-6}
  \Delta w
  =
  2|\nabla_s\xi|^2
  +2|\nabla_t\xi|^2
  +2\langle\xi,\nabla_s\nabla_s\xi+\nabla_t\nabla_t\xi\rangle.
\end{equation}
我们必须证明函数 $w$ 满足引理~\ref{lem:4-3-3} 的假设\sourcefootnote{原书此处写作
“$\Delta w$ 满足引理 4.3.3 的假设”；
但引理~\ref{lem:4-3-3} 的非负性、积分小量与微分不等式均以函数 $w$ 为对象，故改作
“$w$ 满足”。}。由于
\[
  \xi=-J\eta,
  \qquad
  \nabla_s\eta=\nabla_t\xi,
\]
有
\begin{equation}\label{eq:4-3-7}
  \nabla_s\xi+\nabla_t\eta
  =
  \nabla_t(J\xi)-\nabla_s(J\eta)
  =
  (\nabla_\eta J)\xi-(\nabla_\xi J)\eta.
\end{equation}
所以
\[
\begin{aligned}
  \nabla_s\nabla_s\xi+\nabla_t\nabla_t\xi
  &=
  \nabla_s(\nabla_s\xi+\nabla_t\eta)
  +\nabla_t\nabla_s\eta-\nabla_s\nabla_t\eta\\
  &=
  \nabla_s\bigl((\nabla_\eta J)\xi-(\nabla_\xi J)\eta\bigr)
  -R(\xi,\eta)\eta,
\end{aligned}
\]
其中最后一个等号使用了恒等式
\[
  R(\xi,\eta)\eta
  =(\nabla_s\nabla_t-\nabla_t\nabla_s)\eta.
\]
因此
\begin{equation}\label{eq:4-3-8}
  \frac12\Delta w
  =
  |\nabla_s\xi|^2+|\nabla_t\xi|^2
  -\langle R(\xi,\eta)\eta,\xi\rangle
  +\kappa,
\end{equation}
其中误差项 $\kappa$ 为
\[
\begin{aligned}
  \kappa
  &=
  \left\langle
    \xi,
    \nabla_s\bigl((\nabla_\eta J)\xi-(\nabla_\xi J)\eta\bigr)
  \right\rangle\\
  &=
  \left\langle
    \xi,
    (\nabla_\eta J)\nabla_s\xi-(\nabla_\xi J)\nabla_s\eta
  \right\rangle\\
  &\quad+
  \left\langle
    \xi,
    \bigl(\nabla_s(\nabla_\eta J)\bigr)\xi
    -\bigl(\nabla_s(\nabla_\xi J)\bigr)\eta
  \right\rangle.
\end{aligned}
\]
存在常数 $c=c(M,J,g)>0$，使得
\[
  |\nabla_s(\nabla_\eta J)|
  \le c\bigl(|\xi|^2+|\nabla_t\xi|\bigr),
  \qquad
  |\nabla_s(\nabla_\xi J)|
  \le c\bigl(|\xi|^2+|\nabla_s\xi|\bigr).
\]
第一个估计中，$|\xi|^2=|\xi||\eta|$ 的系数涉及 $J$ 的二阶导数估计，而
$|\nabla_t\xi|=|\nabla_s\eta|$ 的系数涉及 $J$ 的一阶导数估计。因此存在常数
$c'=c'(M,J,g)>0$，使得
\[
\begin{aligned}
  \kappa
  &\ge
  -c'|\xi|^4
  -c'|\xi|^2\bigl(|\nabla_s\xi|+|\nabla_t\xi|\bigr)\\
  &\ge
  -\frac12|\nabla_s\xi|^2
  -\frac12|\nabla_t\xi|^2
  -c'(1+c')|\xi|^4.
\end{aligned}
\]
由 \eqref{eq:4-3-8}，存在常数 $c''=c''(M,J,g)>0$，使得
\[
  \Delta w\ge-c''|\xi|^4=-c''w^2.
\]
又因为 $w=\frac12|du|^2$，所以取
$\delta=\pi/(4c'')$，引理~\ref{lem:4-3-3} 即\BookRecordIndex{T0326}推出断言 \textup{(i)}。

为证明 \textup{(ii)}，通过反射把函数
$w:B_{2r}\cap\mathbb H\to\mathbb R$ 延拓到 $B_{2r}$，即对 $t>0$ 定义
\[
  w(s,-t):=w(s,t).
\]
我们证明边界上的法向导数消失：
\[
\begin{aligned}
  \frac12\partial_tw(s,0)
  &=
  \langle\nabla_t\xi,\xi\rangle(s,0)\\
  &=
  \langle\nabla_s\eta,\xi\rangle(s,0)\\
  &=
  \langle\nabla_s(J\xi),\xi\rangle(s,0)\\
  &=
  \langle J\nabla_s\xi,\xi\rangle(s,0)\\
  &=0.
\end{aligned}
\]
倒数第二个等号来自自同态 $\nabla_sJ$ 是\BookRecordIndex{T0350}斜伴随的。最后一个等号则来自：$L$ 是\BookRecordIndex{T0286}全测地
子流形，且 $\xi(s,0)\in T_{u(s,0)}L$，所以
$\nabla_s\xi(s,0)\in T_{u(s,0)}L$。由于 $\partial_tw(s,0)=0$，延拓后的函数 $w$
属于 $C^2$，故由不等式 $\Delta w\ge-c''w^2$ 和引理~\ref{lem:4-3-3} 可得
\textup{(ii)}。多出的因子 $2$ 是因为积分区域为半圆盘。
\end{proof}

\begin{exercise}\label{ex:4-3-5}
设 $S^2$ 是配备面积为 $\pi$ 的\BookRecordIndex{T0363}辛形式的黎曼球面。求使引理~\ref{lem:4-3-1} 中的
平均值估计对映入 $S^2$ 的映射成立时，$\delta$ 的最大值。
\end{exercise}

\begin{remark}\label{rem:4-3-6}
假设 $J$ 与 $\omega$ \BookRecordIndex{T0348}相容。那么每条 \BookRecordIndex{T0003}$J$-全纯曲线 $u:\mathbb C\to M$ 都是\BookRecordIndex{T0105}调和映射，
即
\begin{equation}\label{eq:4-3-9}
  \nabla_s\partial_su+\nabla_t\partial_tu
  =
  (\nabla_{\partial_tu}J)\partial_su
  -(\nabla_{\partial_su}J)\partial_tu
  =0.
\end{equation}
第一个等号来自 \eqref{eq:4-3-7}。为证明 \eqref{eq:4-3-9} 中的第二个等号，使用公式
$\partial_tu=J\partial_su$，以及引理 C.7.1 的 \textup{(C.7.5)} 中
$\omega$-\BookRecordIndex{T0348}相容 $J$ 所满足的恒等式
\[
  (\nabla_{J\partial_su}J)\partial_su
  =
  (\nabla_{\partial_su}J)J\partial_su.
\]
由此，\eqref{eq:4-3-8} 中的函数 $\kappa$ 消失，因而
\begin{equation}\label{eq:4-3-10}
  \frac12\Delta|\partial_su|^2
  =
  |\nabla_s\partial_su|^2
  +|\nabla_t\partial_su|^2
  -\left\langle
    R(\partial_su,\partial_tu)\partial_tu,\partial_su
  \right\rangle.
\end{equation}
若 $(M,\omega,J)$ 的\BookRecordIndex{T0190}截面曲率非正，则右边非负，故函数 $|\partial_su|^2$ 是次调和的。
因此，在这种情形，\BookRecordIndex{T0079}次调和函数的最大值原理（或\BookRecordIndex{T0266}平均值不等式）表明，每条能量有限的
\BookRecordIndex{T0003}$J$-全纯曲线 $u:\mathbb C\to M$ 都是常值映射。换言之，当 $J$ 与 $\omega$ \BookRecordIndex{T0348}相容，且
度量 $\omega(\,\cdot\,,J\,\cdot\,)$ 的\BookRecordIndex{T0190}截面曲率非正时，可以去掉
引理~\ref{lem:4-3-1} 中的小能量假设，因而不存在非常值 $J$-全纯球面。特别地，这对\BookRecordIndex{T0190}截面
曲率非正的 \BookRecordIndex{T0207}凯勒 流形成立。
\end{remark}
